%; whizzy paragraph -pdf xpdf -latex ./whizzypdfptex.sh
%; whizzy-paragraph "^\\\\begin{frame}\\|\\\\emtext"
% latex beamer presentation.
% platex, latex-beamer でコンパイルすることを想定。 

%     Tokyo Debian Meeting resources
%     Copyright (C) 2012 Junichi Uekawa

%     This program is free software; you can redistribute it and/or modify
%     it under the terms of the GNU General Public License as published by
%     the Free Software Foundation; either version 2 of the License, or
%     (at your option) any later version.

%     This program is distributed in the hope that it will be useful,
%     but WITHOUT ANY WARRANTY; without even the implied warreanty of
%     MERCHANTABILITY or FITNESS FOR A PARTICULAR PURPOSE.  See the
%     GNU General Public License for more details.
%     You should have received a copy of the GNU General Public License
%     along with this program; if not, write to the Free Software
%     Foundation, Inc., 51 Franklin St, Fifth Floor, Boston, MA  02110-1301 USA

\documentclass[cjk,dvipdfmx,12pt]{beamer}
\usetheme{Tokyo}
\usepackage{monthlypresentation}

%  preview (shell-command (concat "evince " (replace-regexp-in-string "tex$" "pdf"(buffer-file-name)) "&")) 
%  presentation (shell-command (concat "xpdf -fullscreen " (replace-regexp-in-string "tex$" "pdf"(buffer-file-name)) "&"))
%  presentation (shell-command (concat "evince " (replace-regexp-in-string "tex$" "pdf"(buffer-file-name)) "&"))

%http://www.naney.org/diki/dk/hyperref.html
%日本語EUC系環境の時
\AtBeginDvi{\special{pdf:tounicode EUC-UCS2}}
%シフトJIS系環境の時
%\AtBeginDvi{\special{pdf:tounicode 90ms-RKSJ-UCS2}}

\newenvironment{commandlinesmall}%
{\VerbatimEnvironment
  \begin{Sbox}\begin{minipage}{1.0\hsize}\begin{fontsize}{8}{8} \begin{BVerbatim}}%
{\end{BVerbatim}\end{fontsize}\end{minipage}\end{Sbox}
  \setlength{\fboxsep}{8pt}
% start on a new paragraph

\vspace{6pt}% skip before
\fcolorbox{dancerdarkblue}{dancerlightblue}{\TheSbox}

\vspace{6pt}% skip after
}
%end of commandlinesmall

\title{東京エリアDebian勉強会}
\subtitle{第151回 2017年5月度}
\author{杉本 典充  dictoss@live.jp}
\date{2017年5月20日}
\logo{\includegraphics[width=8cm]{image200607/openlogo-light.eps}}

\begin{document}

\begin{frame}
\titlepage{}
\end{frame}

\begin{frame}{Agenda}
 \begin{minipage}[t]{0.45\hsize}
  \begin{itemize}
  \item 最近あったDebian関連のイベント報告
	\begin{itemize}
	\item 第150回東京エリアDebian勉強会
	\end{itemize}
  \item 事前課題発表
  \end{itemize}
 \end{minipage}
 \begin{minipage}[t]{0.45\hsize}
  \begin{itemize}
   \item Debian Trivia Quiz
  \end{itemize}
 \end{minipage}
\end{frame}

\section{イベント報告}
\emtext{イベント報告}

\begin{frame}{第150回東京エリアDebian勉強会}
\begin{itemize}
\item 2017年4月22日に、IPAさんをお借りして開催しました
\item 参加者は12名でした
\item 「Debianパッケージのカスタマイズについて」 発表者：yy\_y\_ja\_jpさん
\item 同じ会場に集まっていたopenSUSEの方々と交流しました
\end{itemize} 
\end{frame}

\section{事前課題}
\emtext{事前課題}
{\footnotesize
 \begin{prework}{ kenhys }
  \begin{enumerate}
  \item $BH/I=;qNA$r<jD>$7$7$F8x3+$7$F$*$/(B
  \end{enumerate}
\end{prework}

\begin{prework}{ Charles Plessy }
  \begin{enumerate}
  \item r-bioc-* $B%Q%C%1!<%8$N(BQC
  \end{enumerate}
\end{prework}

\begin{prework}{ dictoss }
  \begin{enumerate}
  \item python$B$N(Bsetuptools$B$N3X=,!"(Bpython$B$N(Bdebian$B%Q%C%1!<%8$N:n$jJ}(B
  \end{enumerate}
\end{prework}

\begin{prework}{ yy\_y\_ja\_jp }
  \begin{enumerate}
  \item DDTSS
  \end{enumerate}
\end{prework}

\begin{prework}{ issei }
  \begin{enumerate}
  \item Debconf$B$N%S%G%*4U>^(B
  \end{enumerate}
\end{prework}

\begin{prework}{ Haruki TSURUMOTO(tSU\_RooT) }
  \begin{enumerate}
  \item $BB?J,(BBTS$B$H%a!<%j%9$rFI$`(B
  \end{enumerate}
\end{prework}

\begin{prework}{ ysaito }
  \begin{enumerate}
  \item ($B2sEz$J$7(B)
  \end{enumerate}
\end{prework}

}

\subsection{問題}
%; whizzy-master ../debianmeetingresume201311.tex
% $B0J>e$N@_Dj$r$7$F$$$k$?$a!"$3$N%U%!%$%k$G(B M-x whizzytex $B$9$k$H!"(Bwhizzytex$B$,MxMQ$G$-$^$9!#(B
%

\santaku
{debian$B%W%m%8%'%/%H$N%Q%C%1!<%84XO"$N%5!<%P$N$&$A!"(BFTP$B%5!<%P$NDs6!$r0lItGQ;_$9$k$3$H$,H/I=$5$l$^$7$?!#(BFTP$B%5!<%P$N$&$A!"$=$N8e$b7QB32TF/$9$kM=Dj$N%5!<%P$O$I$l$G$7$g$&$+!#(B}
{ftp://ftp.debian.org}
{ftp://security.debian.org}
{ftp://ftp.upload.debian.org}
{C}
{$B!V(BShutting down public FTP services$B!W$H$$$&I=Bj$G%"%J%&%s%9$5$l$^$7$?!#(Bapt-line$B$K;XDj$9$k%W%m%H%3%k$ODd;_M=Dj$N(B2017-11-01$B$^$G$K(B"http://"$B$KJQ99$7$^$7$g$&!#$J$*!"3+H/<T8~$1$N(B"ftp://ftp.upload.debian.org"$B!"(B"ftp://security-master.debian.org"$B$O$=$N8e$b7QB32TF/$9$k$H$N$3$H$G$9!#(B \url{https://lists.debian.org/debian-announce/2017/msg00001.html}$B!#(B}


\section{セミナー1}
\emtext{Debbugsとのつきあいかた}

\section{Hack time}
\emtext{Hack time}

%\begin{frame}{成果を記入下さい!}
%\end{frame}
  
\section{今後のイベント}
\emtext{今後のイベント}
\begin{frame}{今後のイベント}
  \begin{itemize}
  \item 2017/6/17(土) 東京エリアDebian勉強会（予定）
    \begin{itemize}
    \item 日時、会場、セミナー内容は調整中
    \end{itemize}
  \item 2017/7/15(土) OSC 2017 Hokkaidoへ出張参加
    \begin{itemize}
    \item 出展、またはDebianJPメンバが参加する可能性あり
    \end{itemize}
  \item 2017/7の東京エリアDebian勉強会はお休みとなります
  \item 2017/8/19(土) 東京エリアDebian勉強会（予定）
    \begin{itemize}
    \item 日時、会場、セミナー内容は調整中
    \end{itemize}
  \item 勉強会では発表者を随時募集中です。ぜひともよろしくお願いします。
  \end{itemize}
\end{frame}

\end{document}

;;; Local Variables: ***
;;; outline-regexp: "\\([ 	]*\\\\\\(documentstyle\\|documentclass\\|emtext\\|section\\|begin{frame}\\)\\*?[ 	]*[[{]\\|[]+\\)" ***
;;; End: ***
